\section{Simulation Results}
\label{sec:results}

\subsection{Functional Verification}

We simulate BFS-HBM using Icarus Verilog~\cite{williams2002icarus} with a cycle-accurate
HBM memory model (\texttt{hbm\_mem\_model.sv}) configured for a fixed read latency of
$T_\text{access} = 20$ cycles. The test graph is an 8-vertex undirected graph encoded
in CSR format:

\begin{center}
\texttt{0--1--2--3} \\[2pt]
\texttt{|\ \ \ \ |} \\[2pt]
\texttt{4--5--6--7}
\end{center}

CSR encoding:
\begin{align*}
&\mathtt{row\_ptr} = [0,\,2,\,5,\,8,\,10,\,12,\,15,\,18,\,20] \\
&\mathtt{col\_idx} = [1,4,\; 0,2,5,\; 1,3,6,\; 2,7,\; 0,5, \\
&\phantom{\mathtt{col\_idx} = [}1,4,6,\; 2,5,7,\; 3,6]
\end{align*}

BFS from source vertex 0 at 250~MHz (\texttt{CLK\_PERIOD = 4~ns}).
Table~\ref{tab:sim_results} compares observed and expected BFS levels.

\begin{table}[h]
\centering
\caption{Functional verification: BFS levels, source vertex 0.}
\label{tab:sim_results}
\small
\begin{tabular}{@{}ccrr@{}}
\toprule
Vertex & Exp.\ Level & Obs.\ Level & Discovery (cyc.) \\
\midrule
0 & 0 & 0 & $\sim$11 \\
1 & 1 & 1 & $\sim$68 \\
4 & 1 & 1 & $\sim$74 \\
2 & 2 & 2 & $\sim$136 \\
5 & 2 & 2 & $\sim$142 \\
3 & 3 & 3 & $\sim$264 \\
6 & 3 & 3 & $\sim$270 \\
7 & 4 & 4 & $\sim$396 \\
\bottomrule
\end{tabular}
\end{table}

All 8 vertices are discovered with correct BFS levels. The simulation reports:
\begin{center}
\texttt{BFS complete in 514 cycles. Visited count: 8 / 8.}\\
\texttt{*** ALL 8 VERTICES CORRECT ***}
\end{center}

\subsection{Cycle Budget Breakdown}

Total cycles: 514. We account for the cycle budget analytically using
Equation~(\ref{eq:Tv}) with $T_\text{access} = 20$ and $T_\text{RMW} = 2$:

\begin{table}[h]
\centering
\caption{Cycle budget breakdown by vertex (degree $d_v$).}
\label{tab:cycle_budget}
\small
\begin{tabular}{@{}ccrrrr@{}}
\toprule
Vertex & $d_v$ & $T_\text{ptr}$ & $T_\text{edge}$ & $T_\text{scatter}$ & $T_\text{total}$ \\
\midrule
0 & 2 & 23 & 20 & 4  & 49 \\
1 & 3 & 23 & 20 & 6  & 51 \\
4 & 2 & 23 & 20 & 4  & 49 \\
2 & 3 & 23 & 20 & 6  & 51 \\
5 & 3 & 23 & 20 & 6  & 51 \\
3 & 2 & 23 & 20 & 4  & 49 \\
6 & 3 & 23 & 20 & 6  & 51 \\
7 & 2 & 23 & 20 & 4  & 49 \\
\midrule
Total & 20 & 184 & 160 & 40 & $\approx 400$ + overhead \\
\bottomrule
\end{tabular}
\end{table}

Analytical estimate $\approx 400$ cycles; observed 514 cycles.
The 114-cycle overhead arises from FSM state transition latency, FIFO read handshake
cycles (\textsc{Dequeue}), the Init state bitmap-set latency for the source vertex,
and the final \textsc{CheckEdges} $\to$ \textsc{Dequeue} $\to$ \textsc{Done}
transition sequence when the frontier empties.

\subsection{Utilization}

Useful scatter cycles: $20 \times 2 = 40$ out of 514 total = \textbf{7.8\%}.
This low figure reflects the small test graph: average degree $\bar{d} = 2.5$
makes the per-vertex pointer fetch overhead ($T_\text{ptr} = 23$ cycles) dominate
over the scatter work (5 cycles). From Equation~(\ref{eq:cpe}), utilization at
$\bar{d} = 2.5$ and $T_\text{access} = 20$:

\begin{equation}
\eta_\text{arch}(2.5)
= \frac{2 \times 2.5}{23 + 6.25 + 5} \approx \frac{5}{34.25} \approx 14.6\%
\end{equation}

Utilization rises substantially with graph degree. For $\bar{d} = 40$
(a typical HFT transaction graph):

\begin{equation}
\eta_\text{arch}(40)
= \frac{80}{23 + 100 + 80} \approx \frac{80}{203} \approx 39.4\%
\end{equation}

This converges to the system-level $\eta \approx 40\%$ derived from the HBM
timing model~\cite{fritz2026rmwwall}, confirming that the microarchitecture and
the memory system efficiency bound are consistent.

\subsection{Synthesis}

Yosys synthesis targeting generic gates with reduced parameters
(\texttt{VERTEX\_W=6}, \texttt{FIFO\_DEPTH=16}) reports 847 cells and
a maximum logic depth of 9 levels (excluding BRAM).
The dominant logic is the FSM state register (4 bits), the edge-range registers
(\texttt{edge\_start}, \texttt{edge\_end}, \texttt{edge\_remaining}: $3 \times 32$ bits),
and the beat buffer ($8 \times 32$ bits). Timing closure at 250~MHz on UltraScale+
is expected to be straightforward; the critical path is through the beat buffer
select and bitmap address decode, with no combinational feedback through the
FSM state register.
